% !TeX program = pdflatex
\documentclass[11pt,a4paper]{article}

\usepackage[utf8]{inputenc}
\usepackage[T1]{fontenc}
\usepackage{newtxtext,newtxmath}
\usepackage[margin=2.4cm]{geometry}
\usepackage{amsmath,amssymb,amsthm}
\usepackage{booktabs}
\usepackage{longtable}
\usepackage{array}
\usepackage{enumitem}
\usepackage{xcolor}
\usepackage{csquotes}
\usepackage[colorlinks=true,linkcolor=blue!50!black,citecolor=blue!50!black,
            urlcolor=blue!50!black]{hyperref}
\usepackage{microtype}

% Long monospace identifiers (copula::serialIndepTestSim etc.) cannot hyphenate;
% give TeX slack rather than letting them punch through the margin.
\emergencystretch=3em
\hbadness=2000

\newcommand{\Vtilde}{\widetilde{V}^{0}_{\tau,\hat\gamma}}
\newcommand{\Wbridge}{W^{0}}
\newcommand{\gh}{\hat\gamma}
\newcommand{\Ntau}{N(\tau)}

% Provenance tags
\newcommand{\PROVED}{\textsc{\textcolor{blue!55!black}{proved}}}
\newcommand{\SIMULATED}{\textsc{\textcolor{orange!70!black}{simulated}}}
\newcommand{\ASSERTED}{\textsc{\textcolor{purple!70!black}{asserted}}}
\newcommand{\MEASURED}{\textsc{\textcolor{green!45!black}{measured}}}

\title{Statistical Instruments of the QRE Licensing Battery\\
\large Definitions, provenance, calibration, and implementation status}
\author{Reference document for \texttt{qretool}}
\date{}

\begin{document}
\maketitle

\begin{abstract}
\noindent
This document defines every statistical instrument used or considered in the licensing
stage of the Quantum Reliability Engineering framework. For each instrument it states
the object being computed, the null hypothesis, the alternative it is directed at, the
calibration mechanism and why that mechanism rather than another, the censoring
behaviour, the published provenance down to equation number, and whether the
implementation imports a library or transcribes an equation. Load-bearing claims carry
a provenance tag: \PROVED{} for a theorem or exact algebraic result, \SIMULATED{} for a
Monte Carlo finding, \ASSERTED{} for a textbook or review statement given without proof
in the cited passage, \MEASURED{} for an empirical data result.
\end{abstract}

\tableofcontents

% ==============================================================
\section{Three worlds, and which one this battery lives in}
% ==============================================================

Statistical questions about the window durations fall into three groups that need
different instruments and different prerequisites. Mixing them is the single most common
source of confusion in this project, so they are separated first.

\begin{center}
\small
\begin{tabular}{@{}p{0.16\textwidth}p{0.36\textwidth}p{0.24\textwidth}p{0.18\textwidth}@{}}
\toprule
\textbf{World} & \textbf{Question} & \textbf{Instruments} & \textbf{Prerequisite}\\
\midrule
A. Licensing &
May serially observed windows be treated as parallel replicate lifetimes? &
C1, C2, C3, C5, C6 &
window table\\[2pt]
B. Description &
What shape is the lifetime distribution? Is an exponential model defensible? &
Weibull shape CI, TTT plot, probability plots, censored GoF &
Kaplan-Meier\\[2pt]
C. Model adequacy &
Is the fitted regression structure correct? &
Lin-Wei-Ying, Schoenfeld, PH tests &
a fitted regression\\
\bottomrule
\end{tabular}
\end{center}

\medskip
\noindent
This battery is World A only. World B is deferred until the product-limit estimator
lands. World C is gated on a model-fit branch that does not exist, which is why the
Lin-Wei-Ying check is documented in Section~\ref{sec:lwy} and not implemented.

\subsection{The one assumption behind the whole battery}

The estimator, the confidence band, and the $k$-sample log-rank test all descend from
the Aalen multiplicative intensity model (Andersen, Borgan, Gill and Keiding, 1993,
\S IV.3 and \S V.2), in which the intensity of the counting process factorises as
$\lambda(t) = \alpha(t)\,Y(t)$ with $Y$ predictable. \PROVED

For a \emph{single} unit observed sequentially, with the at-risk process resetting at
each renewal and the baseline hazard a deterministic function of backward recurrence
time, that model is equivalent to independent and identically distributed window
durations, which is the definition of a renewal process. \PROVED

Two consequences follow and they fall together. The Greenwood band and the log-rank
merge verdict are not two separate tools resting on two separate assumptions; they are
two consequences of one assumption. A check failure revokes both, never one alone.

Hence the entire licensing question reduces to: \emph{do the carved windows form a
renewal sequence, or at least a stationary weakly dependent one?}

\subsection{What counts as a \enquote{process}, and what \texorpdfstring{$m$}{m} is}
\label{sec:whatism}

This is the single most important mapping in the document and getting it backwards makes
the whole battery look degenerate.

In Kvaløy and Lindqvist a \emph{process} is one observation unit: a counting process
$N_j(t)$ watched from $0$ to its own censoring time $\tau_j$, inside which a random
number $N_j(\tau_j)$ of \emph{events} occur with \emph{gaps} $X_i$ between them.

The mapping used here is:

\begin{center}
\begin{tabular}{@{}ll@{}}
\toprule
\textbf{Their object} & \textbf{Our object}\\
\midrule
process $j$ & one gap-free stretch of the scan record\\
event $T_i$ & a down-crossing, i.e.\ a window death\\
gap $X_i$ & one in-spec window duration\\
$\tau_j$ & when watching stopped for that stretch\\
$N_j(\tau_j)$ & the number of complete windows in that stretch\\
\bottomrule
\end{tabular}
\end{center}

\medskip
\noindent
\textbf{So $m=1$ does not mean one window.} It means one uninterrupted record containing
$N$ windows, and $N$ is in the tens to low hundreds for this data. A single-process
analysis at $N=355$ is the ordinary case, not a degenerate one.

$m>1$ arises when the record is interrupted, because events cannot be observed during a
read gap and the two sides cannot be joined without fabricating an inter-event gap
across it. So $m$ is one more than the number of gaps that split the record, not the
number of windows. Gaps are infrequent, so $m$ is usually small and often $1$.

\subsection{Which tests survive at \texorpdfstring{$m>1$}{m>1}}
\label{sec:mgt1}

The paper's \S 4.2 is explicit about which extensions are easy and which are not, and the
distinction is purely one of algebra: a functional that is \emph{linear} in the bridge
carries through the weighted sum $\sum_j A_j \Vtilde[j]$, and a functional that squares
or takes a supremum does not.

\begin{center}
\small
\begin{tabular}{@{}p{0.13\textwidth}p{0.10\textwidth}p{0.30\textwidth}p{0.38\textwidth}@{}}
\toprule
\textbf{Test} & \textbf{Linear?} & \textbf{$m>1$ route} & \textbf{Paper's verdict}\\
\midrule
LR & yes & closed form, eq.~(16) & recommended for monotone trend at any $m$\\
ELR & yes & closed form, by the same linearity & \enquote{This also applies to the extended Lewis-Robinson test}\\
KS & no & sum of statistics & not pursued\\
CvM & no & sum of statistics, normal approximation & \enquote{works fairly well}\\
AD & no & sum of statistics, normal approximation & \enquote{less well \dots{} due to the very skew distribution}\\
\bottomrule
\end{tabular}
\end{center}

\medskip
\noindent
Verbatim from \S 4.2 on the linear cases: \enquote{This corresponds to what we did for the
Lewis-Robinson test in the previous subsection, but here things are easy due to the
linearity of integrals. This also applies to the extended Lewis-Robinson test.} And on
the rest: \enquote{For the remaining tests it is not straightforward, however, to derive
explicit expressions for the test statistics, and it is neither clear what would be the
best weights to use.} \ASSERTED

\paragraph{Is $m>1$ with a non-monotonic alternative uncovered?} No, but only because of
a choice the implementation makes that the paper does not. Three routes exist:

\begin{enumerate}[leftmargin=1.6em,itemsep=3pt]
\item \textbf{ELR at $m>1$.} Closed form by linearity, and ELR is directed at
non-monotonic trend. Cost: $a$ must be prespecified (Section~\ref{sec:missing}).
\item \textbf{CvM at $m>1$.} Sum of per-process statistics with a normal approximation,
which the authors report works fairly well. This is the paper's own recommendation for
the regime, and it is the strongest argument for implementing CvM.
\item \textbf{AD at $m>1$ by permutation.} The route this implementation takes. The
statistic is the unweighted sum of per-segment eq.~(7) values, calibrated by
within-segment permutation. \emph{This sidesteps the paper's problem rather than solving
it}: the objection in \S 4.2 is to the normal \emph{approximation} to the sum, and a
permutation calibration makes no distributional claim at all. It is exact under
independent and identically distributed gaps.
\end{enumerate}

\noindent
Route 3 is defensible and arguably better than routes 1 and 2, since it needs neither a
tuning parameter nor an approximation the authors themselves distrust. But it is an
extension beyond the source, so it carries no published size evidence and \textbf{must be
validated in the bench before any $m>1$ Anderson-Darling verdict is reported}. Until then
the honest statement is that the $m>1$ non-monotonic cell rests on the implementation's
own calibration, not on the paper's.

\subsection{Vocabulary: parametric, nonparametric, semiparametric}
\label{sec:semiparametric}

Reliability engineering canon is parametric-first, biostatistics canon is
semiparametric-first, and the vocabulary gap is a recurring source of trouble. There are
three tiers, not two.

\begin{description}[leftmargin=1.6em,itemsep=2pt]
\item[Parametric.] The whole distribution is written down. Weibull, exponential,
lognormal. A finite parameter vector determines everything.
\item[Nonparametric.] Nothing is written down. Kaplan-Meier, Nelson-Aalen. The
distribution is an infinite-dimensional unknown estimated directly.
\item[Semiparametric.] Part is written down, the rest is left free. Cox is canonical:
$\lambda(t \mid Z) = \lambda_0(t)\exp(\beta' Z)$, where the covariate effect is
parametric with finite $\beta$ and the baseline $\lambda_0(t)$ is left completely
unspecified. Aalen's additive model is the same idea with a different link.
\end{description}

This tier is why Lin-Wei-Ying does not belong in the battery. It checks a fitted
semiparametric \emph{regression structure}, not a distribution family and not
independence. See Section~\ref{sec:lwy}.

% ==============================================================
\section{The construction shared by C1 and C2}
% ==============================================================

C1 and C2 are not two unrelated tests. They are two functionals of one object, and
understanding the object explains both at once, explains why the paper offers four
tests rather than two, and explains what calibration is available.

\subsection{The functional central limit theorem}

Let a renewal process be observed from $t = 0$ with event times $T_1, T_2, \dots$ and
gap times $X_i = T_i - T_{i-1}$, independent and identically distributed with
$E(X_i) = \mu$, $\mathrm{Var}(X_i) = \sigma^2 < \infty$, and coefficient of variation
$\gamma = \sigma/\mu$. Let $N(t)$ count events in $(0,t]$.

Billingsley (1999, Theorem 14.6) gives, for
$V_{t,\mu,\sigma}(s) = \mu^{3/2}\bigl(N(st) - st/\mu\bigr)/\bigl(\sigma\sqrt{t}\bigr)$,
weak convergence $V_{t,\mu,\sigma} \Rightarrow W$ to the Wiener measure. \PROVED

Tying the process down at both ends gives Kvaløy and Lindqvist's Theorem 1: with
$V^0_{t,\mu,\sigma}(s) = V_{t,\mu,\sigma}(s) - s\,V_{t,\mu,\sigma}(1)$, one has
$V^0_{t,\mu,\sigma} \Rightarrow \Wbridge$, a Brownian bridge. Their Corollary 1 replaces
the unknown parameters by the coefficient of variation alone:
%
\begin{equation}
\Vtilde(s) \;=\; \frac{1}{\gh}\,\frac{N(s\tau) - s\,\Ntau}{\sqrt{\Ntau}}
\;\Longrightarrow\; \Wbridge
\qquad\text{as } \tau \to \infty .
\label{eq:corollary1}
\end{equation}
\PROVED{} (Kvaløy and Lindqvist, arXiv:1802.08339, eq.~(3) and Corollary 1.)

\subsection{What this buys}

Equation~\eqref{eq:corollary1} says: under the renewal null, the observed counting
process, normalised and tied down, looks like a Brownian bridge. Any functional that
measures deviation from a Brownian bridge is therefore a trend test, and its null
distribution is the distribution of that functional applied to $\Wbridge$, which is
known classically.

This is why the paper produces a \emph{class}. Four standard functionals give four
tests:

\begin{center}
\small
\begin{tabular}{@{}llll@{}}
\toprule
\textbf{Functional of $\Wbridge$} & \textbf{Test} & \textbf{Eq.} & \textbf{Limit law}\\
\midrule
$\int_0^1 \Wbridge(s)\,ds$ & Lewis-Robinson & (4) & $N(0, 1/12)$, scaled to $N(0,1)$\\
$\sup_s |\Wbridge(s)|$ & Kolmogorov-Smirnov & (5) & Kolmogorov distribution\\
$\int_0^1 \Wbridge(s)^2\,ds$ & Cramér-von Mises & (6) & classical CvM\\
$\int_0^1 \Wbridge(s)^2 / \bigl(s(1-s)\bigr)\,ds$ & Anderson-Darling & (7) & classical AD\\
\midrule
$\int_0^a \Wbridge - \int_a^1 \Wbridge$ & Extended Lewis-Robinson & (8), (9) &
$N\bigl(0, \tfrac{1}{12} - a^2(1-a)^2\bigr)$\\
\bottomrule
\end{tabular}
\end{center}

\medskip
\noindent
The role of $\gh$ deserves emphasis because it is the whole difference between this
family and the classical Poisson-based tests. Dividing by the estimated coefficient of
variation removes the renewal distribution's dispersion. Without it, eq.~(4) is exactly
the Laplace statistic, whose null is a homogeneous Poisson process and which
over-rejects under overdispersion. That over-rejection is the original observation of
Lewis and Robinson (1974). \PROVED

\subsection{Estimating \texorpdfstring{$\gamma$}{gamma}}

Three estimators are given in the paper's \S 3.6 and justified in its Appendix 1.

\begin{enumerate}[leftmargin=1.6em,itemsep=2pt]
\item Sample mean and standard deviation of the \emph{completely observed} gaps,
$\gh = \hat\sigma / \hat\mu$. Consistent as $\tau\to\infty$, but discards the censored
final gap $\tau - T_{\Ntau}$.
\item Estimators using the censored time, their eq.~(10):
$\tilde\mu = \tau/\Ntau$ and
$\tilde\sigma^2 = \Ntau^{-1}\bigl\{\sum X_i^2 + (\tau - T_{\Ntau})^2\bigr\} - \tilde\mu^2$.
\item The successive-difference estimator, their eq.~(11):
$\sigma^{*2} = \bigl(2(\Ntau-1)\bigr)^{-1}\sum (X_{i+1}-X_i)^2$.
\end{enumerate}

The paper explicitly declines the third: it tends to be smaller under positive
dependence between subsequent gaps, which inflates the test statistic and raises
rejection, and they report \enquote{apparent less satisfactory significance level
properties} in their simulations. \SIMULATED

\medskip
\noindent\textbf{This choice matters here more than it does in the paper.} Estimator 3
is deliberately sensitive to exactly the neighbour dependence C5 and C6 are built to
detect. Using it would entangle the trend test with the dependence test. The
implementation exposes estimators 1 and 2 as \texttt{GAMMA\_COMPLETE} and
\texttt{GAMMA\_TRUNCATED}, and does not offer estimator 3. That is the right call and
should be stated when the choice is questioned.

% ==============================================================
\section{C1 -- Lewis-Robinson, time censored}
% ==============================================================

\paragraph{Object.} The signed area under the tied-down normalised counting process,
scaled to be asymptotically standard normal:
%
\begin{equation}
\mathrm{LR} \;=\; -\sqrt{12}\int_0^1 \Vtilde(s)\,ds
\;=\; \frac{1}{\gh}\cdot\frac{\sqrt{12}}{\tau\sqrt{\Ntau}}
\left[\sum_{i=1}^{\Ntau} T_i - \frac{\Ntau}{2}\,\tau\right].
\label{eq:lr}
\end{equation}
\PROVED{} (their eq.~(4)).

\paragraph{Reading it.} $\sum_i T_i$ is the total of the event times, and
$(\Ntau/2)\tau$ is its expectation when events are uniform on $[0,\tau]$, which is what
no trend means. So LR is a centred, scaled \emph{are the events early or late}
contrast. Positive means increasing intensity, negative means decreasing. It is used
two-sided here.

\paragraph{Null.} Renewal process: gaps independent and identically distributed.

\paragraph{Directed at.} Monotone trend. \enquote{The resulting test will primarily have
power against deviations from an RP caused by monotonic trends.} \ASSERTED

\paragraph{Multi-process form, and why it is the one implemented.} With $m$ independent
processes, weights $A_j \propto \gh_j \tau_j \sqrt{N_j(\tau_j)}$ (their eq.~(14),
justified by the limit in eq.~(15)) inserted into eq.~(13) give
%
\begin{equation}
\mathrm{LR}^m \;=\;
\frac{\sqrt{12}}{\sqrt{\sum_{k=1}^{m}\gh_k^2\,\tau_k^2\,N_k(\tau_k)}}
\sum_{j=1}^{m}\left[\sum_{i=1}^{N_j(\tau_j)} T_{ij} - \frac{N_j(\tau_j)}{2}\,\tau_j\right].
\label{eq:lrm}
\end{equation}
\PROVED{} (their eq.~(16)).

Setting $m = 1$ in \eqref{eq:lrm} gives
$\sqrt{12}\,\bigl[\textstyle\sum T_i - N\tau/2\bigr] \big/ \sqrt{\gh^2\tau^2 N}
= (1/\gh)\cdot\sqrt{12}/(\tau\sqrt{N})\cdot[\,\cdot\,]$, which is \eqref{eq:lr}
exactly. So a separate single-process code path would be a second copy of one formula,
free to drift. Only \eqref{eq:lrm} is implemented and the collapse is asserted
numerically in the test suite rather than trusted to the algebra above.

\paragraph{Why $m \neq 1$ is not hypothetical here.} A read gap means events cannot be
observed during it, so the record splits into gap-free segments, each with its own
$\tau_j$. The multi-process form is the gap case, not an optional generalisation. On a
record with three gaps, $m = 4$.

\paragraph{Calibration.}

\emph{Asymptotic, $N(0,1)$ two-sided.} The chain is: Billingsley (1999, Theorem 14.6)
gives the functional central limit theorem for renewal processes; Kvaløy and Lindqvist's
Theorem 1 and Corollary 1 tie it down to a Brownian bridge, using Billingsley (1999,
Theorem 3.1), the converging-together lemma, to replace $\mu,\sigma$ by $\gh$; their
\S 3.1 then states that $\int_0^1 \Wbridge(s)\,ds$ \enquote{is normally distributed with
expectation 0 and variance $1/12$}, so the $-\sqrt{12}$ scaling in eq.~(4) gives a
standard normal limit. For $m>1$ the additional ingredient is their Lemma 1, that a
weighted sum $\sum_j a_j W_j^0$ of independent Brownian bridges with $\sum a_j^2 = 1$ is
itself a Brownian bridge, plus the weight limit in their eq.~(15) establishing
$A_j \xrightarrow{p} a_j$. \PROVED{} (Kvaløy and Lindqvist 2020, Theorem 1, Corollary 1,
Lemma 1, \S\S 3.1 and 4.1; Billingsley 1999, Theorems 14.6 and 3.1.)

\emph{Permutation of the completely observed gaps within each segment.} Two separate
claims, with separate provenance:

\begin{enumerate}[leftmargin=1.6em,itemsep=2pt]
\item \emph{The authors sanction it.} Their \S 7: the asymptotic critical values
\enquote{are hence only approximate when used in small and medium sized samples}, and
\enquote{in such cases an alternative procedure would be to simulate the null
distribution of the test statistic by a permutation approach, permuting the order of the
completely observed interevent times}. They cite Lawless, Çiğşar and Cook
(\emph{Technometrics} 54:147--158, 2012) for its validity under \emph{time} censoring,
which is the non-obvious part, and confirm it in unreported simulations.
\ASSERTED{} by the authors, \SIMULATED{} by Lawless et al.
\item \emph{It is exact here, not merely valid.} Under the renewal null the gaps within a
segment are exchangeable, so the randomisation hypothesis holds and a permutation test
has exact level $\alpha$ (Lehmann and Romano, \emph{Testing Statistical Hypotheses},
\S 15.2; Romano and Tirlea 2022, \S\S 2--3). What makes it exact for \emph{this}
statistic specifically is an invariance that must be checked and is easy to check:
$\tau_j$, $N_j$ and $\gh_j$ are all invariant to reordering the gaps inside a segment,
because each is a symmetric function of the gap multiset. The only quantity a permutation
moves is $\sum_i T_{ij}$. Consequently the permuted statistic has the same denominator as
the observed one, and the comparison is a like-for-like reordering of one numerator.
\PROVED
\end{enumerate}

\noindent
The practical value of that exactness is measurement, not safety: because the permutation
$p$-value is exact, the asymptotic-versus-permutation gap in the bench is a clean reading
of \emph{asymptotic error}, rather than a comparison of two approximations whose
disagreement could be blamed on either.

\paragraph{Censoring.} Native. The incomplete final gap $\tau - T_{\Ntau}$ is retained
by construction; it is the whole point of the time-censored formulation. \PROVED

% ==============================================================
\section{C2 -- Anderson-Darling type, time censored}
% ==============================================================

\paragraph{Object.}
$\mathrm{AD} = \int_0^1 \Vtilde(s)^2/\bigl(s(1-s)\bigr)\,ds$, which by
\eqref{eq:corollary1} converges to $\int_0^1 \Wbridge(s)^2/(s(1-s))\,ds$, the classical
Anderson-Darling limit (Anderson and Darling, 1952, 1954). Evaluated,
%
\begin{equation}
\begin{split}
\mathrm{AD} = \frac{1}{\gh^{2}}\,\frac{1}{\Ntau}
\Biggl\{&\sum_{i=1}^{\Ntau-1}\left[(\Ntau-i)^2\ln\frac{\tau - T_i}{\tau - T_{i+1}}
+ i^2 \ln\frac{T_{i+1}}{T_i}\right]\\
&+ \Ntau^{2}\left[\ln\frac{\tau}{\tau - T_1} + \ln\frac{\tau}{T_{\Ntau}} - 1\right]\Biggr\}.
\end{split}
\label{eq:ad}
\end{equation}
\PROVED{} (their eq.~(7)).

\paragraph{Reading it.} With $u_i = T_i/\tau$, eq.~\eqref{eq:ad} at $\gh = 1$ is exactly
the classical Anderson-Darling statistic testing the event times for uniformity on
$[0,\tau]$. The $1/\gh^2$ generalises it from a Poisson null to a renewal null. The
weight $1/(s(1-s))$ is what distinguishes it from Cramér-von Mises: it puts more weight
on the beginning and end of the observation interval.

\paragraph{Null.} Renewal process.

\paragraph{Directed at.} Monotonic \emph{and} non-monotonic trend. Because the
integrand is squared, deviation in either direction contributes. This is why the paper
recommends it as the general-purpose single-process test: \enquote{the Anderson-Darling
test turns out to have attractive properties when used as a test for general
alternatives, both monotonic and non-monotonic trends.} \ASSERTED

\paragraph{On the critical value $2.492$.}
The rejection rule \enquote{reject at $\alpha = 0.05$ when $\mathrm{AD} > 2.492$} is
correct and is the classical asymptotic Anderson-Darling critical value. Nothing is
being rederived. The implementation evaluates the same limiting law as a continuous
cumulative distribution function, using the \texttt{adinf} approximation of Marsaglia
and Marsaglia (\emph{Journal of Statistical Software} 9(2), 2004), quoted to
$|\text{error}| < 2\times 10^{-6}$, and it returns $0.05001$ at $z = 2.492$ and
$0.10001$ at $z = 1.933$. Two reasons for the CDF rather than the table:
\begin{enumerate}[leftmargin=1.6em,itemsep=1pt]
\item A $p$-value is strictly more informative than a binary verdict, and the ledger
records $p$-values so that a near-miss is visible.
\item The bench needs a continuous $p$-value to measure empirical size at all, since
empirical size is the rejection rate computed from $p$-values.
\end{enumerate}
Marsaglia and Marsaglia's \emph{finite-$N$} correction \texttt{errfix} is deliberately
not applied: it is fitted for the Anderson-Darling statistic of $N$ independent uniform
variables, whereas \eqref{eq:ad} is additionally divided by an \emph{estimated}
$\gh^2$, so the finite-$N$ null here is not the finite-$N$ Anderson-Darling null.
Applying that correction would be a precision claim the derivation does not support.

\paragraph{Where the recalibration mandate comes from.}
It is the authors', in their \S 7:
\enquote{The derived test statistics are based on asymptotic results for renewal
processes. The calculated critical values are hence only approximate when used in small
and medium sized samples.} They then name the fallback:
\enquote{simulate the null distribution of the test statistic by a permutation approach,
permuting the order of the completely observed interevent times}, citing Lawless,
Çiğşar and Cook (\emph{Technometrics} 54, 2012) for its validity under time censoring
and confirming it in unreported simulations. \ASSERTED{} by the authors,
\SIMULATED{} by Lawless et al.

So the permutation calibration in the implementation is not an invention. It is the
paper's own prescription for the regime this project sits in.

\paragraph{Multi-process.} See Section~\ref{sec:mgt1} for the full picture. In short:
there is no asymptotic calibration for $m > 1$, and the implementation raises rather than
approximating. The paper's \S 4.2 states the normal
approximation to a weighted sum \enquote{works fairly well for the Cramér-von Mises
test, but less well for the Anderson-Darling test due to the very skew distribution of
the Anderson-Darling statistic}, and its \S 5 drops Anderson-Darling for $m > 1$ in
favour of Cramér-von Mises. The multi-process statistic implemented is the unweighted
sum of per-segment values, calibrated by within-segment permutation, which needs no
asymptotic distribution and is exact under independent and identically distributed
gaps. \ASSERTED

\paragraph{Censoring.} Native, same construction as C1.

% ==============================================================
\section{Two tests from the same paper that are not implemented}
\label{sec:missing}
% ==============================================================

This section exists because the omissions are real and should be decided rather than
inherited.

\subsection{Extended Lewis-Robinson, eqs.~(8) and (9)}

\paragraph{This is a different test, not a variant of eq.~(4).} Equation~(8) is
%
\begin{equation}
\int_0^a \Vtilde(s)\,ds \;-\; \int_a^1 \Vtilde(s)\,ds ,
\qquad 0 \le a \le 1,
\label{eq:elr8}
\end{equation}
%
a functional that splits the bridge at $a$ and gives the two halves opposite signs. It
is not obtained by specialising eq.~(16); eq.~(16) is the $m$-process generalisation of
eq.~(4), and eq.~\eqref{eq:elr8} is a different functional of the same bridge. The
scaled statistic is their eq.~(9),
%
\begin{equation}
\mathrm{ELR} = \frac{1}{\gh}\cdot
\frac{1}{\tau\sqrt{\Ntau}\sqrt{\tfrac{1}{12} - a^2(1-a)^2}}
\left\{\sum_{i=1}^{\Ntau}\bigl|T_i - a\tau\bigr|
- \Bigl(\tfrac{1}{2} - a(1-a)\Bigr)\tau\,\Ntau\right\}.
\label{eq:elr9}
\end{equation}
\PROVED

\paragraph{What Appendix 2 actually proves, and its cost.}
Appendix 2 derives \eqref{eq:elr9} from \eqref{eq:elr8} by integrating the indicator
representation $N(t) = \sum_i I_{[T_i,\tau)}(t)$ over $[0,a\tau]$ and $[a\tau,\tau]$
separately and subtracting, and then proves the null law. Writing
$U = \int_0^a \Wbridge$ and $V = \int_a^1 \Wbridge$, and using
$k(s,t) = \min(s,t) - st$:
$\mathrm{Var}(U) = a^3(4-3a)/12$ and
$\mathrm{Var}(V) = 1/12 - a^2(6-8a+3a^2)/12$. The covariance is avoided by the identity
$\mathrm{Var}(U+V) = 1/12$, giving
$\mathrm{Var}(U-V) = 2(\mathrm{Var}\,U + \mathrm{Var}\,V) - \mathrm{Var}(U+V)
= 1/12 - a^2(1-a)^2$. \PROVED

Three costs follow, all of them stated by the authors or immediate from the algebra:

\begin{enumerate}[leftmargin=1.6em,itemsep=3pt]
\item \textbf{$a$ must be fixed in advance.} \enquote{A disadvantage of the above test
is that the value of $a$ has to be given. One possibility would of course be to allow an
adaptive choice of $a$. This will, however, destroy the above distributional properties,
and we will therefore not pursue this approach here.} \ASSERTED{} This is the binding
constraint. Choosing $a$ by looking at the data invalidates \eqref{eq:elr9}.

\item \textbf{At $a = 1/2$ it is tuned to a symmetric bathtub.} The variance falls from
$1/12$ to $1/12 - 1/16 = 1/48$, and the authors note the test is
\enquote{particularly well suited for picking up non-monotonic trends which are
symmetric around the mid-point of the observation interval}. A non-monotonic trend whose
turning point is not near $\tau/2$ is not what this test is aimed at.

\item \textbf{It is worse than LR against monotone trend.} The paper reports it only for
non-monotonic trend \enquote{as it has inferior power against monotonic trends}.
\SIMULATED
\end{enumerate}

\paragraph{Verdict for this project.} Anderson-Darling already covers non-monotonic
trend without requiring a choice of $a$, and their Figure 4 shows both detecting the
bathtub. The Extended Lewis-Robinson buys a sharper instrument aimed at one specific
shape at the cost of a prespecified tuning parameter. Recommendation: implement it only
if a physical argument fixes $a$ a priori. There is no such argument at present, so the
honest position is that the test is documented, understood, and declined on the
tuning-parameter objection, which is the same objection already used to retire the block
bootstrap.

\subsection{Cramér-von Mises, eq.~(6)}

\paragraph{This omission is harder to defend and should be closed.}
%
\begin{equation}
\begin{split}
\mathrm{CvM} = \frac{1}{\gh^2}\frac{1}{\Ntau}
\Biggl\{&\sum_{i=0}^{\Ntau-1}\left[i^2\frac{X_{i+1}}{\tau}
- i\,\Ntau\,\frac{T_{i+1}^2 - T_i^2}{\tau}\right]\\
&+ \Ntau^2\left[\frac{T_{\Ntau}^2}{\tau^2} - \frac{T_{\Ntau}}{\tau}
+ \frac{1}{3}\right]\Biggr\}.
\end{split}
\label{eq:cvm}
\end{equation}
\PROVED

Same null, same sensitivity to monotone and non-monotone trend as Anderson-Darling,
differing only in the weight function. For a single process the paper drops it, since
\enquote{this test had less power than the Anderson-Darling test}. \SIMULATED

\medskip
\noindent\textbf{But for $m > 1$ the paper reverses the preference}: \enquote{for
several processes we do not include the Anderson-Darling test as the Cramér-von Mises
test had better level properties in this case}. \SIMULATED

Since $m > 1$ is the gap case and gaps are routine in this data, Cramér-von Mises is the
paper's own recommendation for a regime the battery actually operates in. It is
approximately thirty lines from \eqref{eq:cvm}, uses the same $\gh$ machinery and the
same permutation harness already built, and closes a gap the source paper names
explicitly. \textbf{Recommendation: implement.}

\subsection{Kolmogorov-Smirnov, eq.~(5)}

Dominated and correctly dropped. The paper: \enquote{The Kolmogorov-Smirnov test is less
powerful than the other tests}, and it is \enquote{too conservative in the overdispersed
case}. \SIMULATED{} Overdispersion is the expected regime for threshold-excursion
durations, so the specific failure mode is the one that would bite. No action.

\subsection{On rank tests in reliability}

A note on the concern that rank statistics are a time-series import with no standing in
recurrent-event analysis. They have standing. Kvaløy and Lindqvist cite Mann (1945),
\emph{Nonparametric tests against trend}, \emph{Econometrica} 13:245--259, as the first
nonparametric test under the renewal null, and it operates on reverse arrangements of
the gap sequence. Rank-based inference on interevent times is native to the field.

The distinction that matters is the \emph{alternative}, not the machinery. Mann's rank
test is directed at monotone \emph{trend}, which makes it redundant with Lewis-Robinson:
same null, same directed alternative, high correlation in practice. C5 and C6 use rank
machinery but are directed at \emph{neighbour dependence}, which no trend test detects.
That is why they are not redundant with C1 and C2, and it is also why Mann is not in the
battery.

% ==============================================================
\section{C3 -- Empirical copula serial independence}
% ==============================================================

\paragraph{Object.} A Cramér-von Mises functional of the empirical copula process of the
lagged vector $(X_i, X_{i+1}, \dots, X_{i+p})$, decomposed by Möbius inversion into
one statistic per subset of lags, so that the dependogram shows which lags carry the
dependence.

\paragraph{Provenance.} Ghoudi, Kulperger and Rémillard, \emph{Journal of Multivariate
Analysis} 79(2):191--218 (2001); Genest and Rémillard, \emph{TEST} 13(2):335--369
(2004); Kojadinovic and Yan, \emph{Annals of the Institute of Statistical Mathematics}
63(2):347--373 (2011). \PROVED

\paragraph{Null.} Full serial independence, jointly across all lags considered. This is
strictly stronger than the zero-rank-autocorrelation null of C5 and C6, which is the
reason it is in the battery despite the R dependency.

\paragraph{Directed at.} Any pairwise and higher-order serial dependence in the copula,
including dependence with zero rank correlation.

\paragraph{Calibration.} The univariate statistic is asymptotically distribution-free,
so critical values come from Monte Carlo simulation of the limiting independent null,
which is what \texttt{copula::serialIndepTestSim} does before
\texttt{copula::serialIndepTest} applies it. The multivariate Kojadinovic-Yan form is
\emph{not} distribution-free and the package documentation states that critical values
and $p$-values come from a bootstrap or permutation of the pseudo-observations. \PROVED

\paragraph{Censoring.} None. No right-censored variant of this statistic has been
located, so censored windows are dropped. Note that \enquote{drop censored windows} is
an inference from the absence of a censored variant, not a sourced rule, and the ledger
must report the dropped count for this reason.

\paragraph{Implementation status.} This is the one check that imports rather than
transcribes, over a file bridge to \texttt{Rscript}. \textbf{As of this writing R is
absent from the development machine, so C3 has never executed and carries no bench
evidence.} Its silence in the ledger is an environment fact, not a pass.

\paragraph{Ties caveat.} The distribution-free property is stated for continuous
observations. Excursion durations carved on a read grid are integer multiples of the
read spacing and carry heavy ties at loose thresholds, which violates the assumption the
calibration rests on. This is a first-class factor in the bench, not a footnote.

% ==============================================================
\section{C5 -- Studentized max-lag rank autocorrelation}
% ==============================================================

\paragraph{Object.}
$T = \max_{h=1..H}\bigl| (r_h - \mu_h)/s_h \bigr|$, where $r_h$ is the within-segment
rank autocorrelation of the duration sequence at lag $h$ and $(\mu_h, s_h)$ are that
lag's moments across the augmented permutation ensemble.

\paragraph{Null.} Exchangeability of the duration sequence, hence serial independence
under the renewal null.

\paragraph{Directed at.} Localised neighbour dependence. If the durations carry a lag-2
relation and nothing else, the maximum picks it up at full strength where a portmanteau
statistic dilutes it across every lag examined.

\paragraph{Provenance of the statistic.} There is no single paper to cite for
\enquote{max-over-lags studentized rank autocorrelation with segmented permutation},
because it is assembled from three pieces, each of which does have a source:

\begin{enumerate}[leftmargin=1.6em,itemsep=2pt]
\item \emph{Rank autocorrelation as a randomness statistic.} The rank version of von
Neumann's ratio is Bartels, R.\ (1982), \enquote{The Rank Version of von Neumann's Ratio
Test for Randomness}, \emph{JASA} 77(377):40--46, implemented as
\texttt{randtests::bartels.rank.test} and \texttt{lawstat::bartels.test}. Bartels'
statistic is a lag-1 rank object and his own Monte Carlo shows it has far greater power
than the runs-up-and-down test. Our $r_1$ is the Spearman form of the same idea, so
\texttt{bartels.rank.test} is a legitimate external cross-check on the lag-1
\emph{statistic}, though not on the max-over-lags \emph{test}. \SIMULATED
\item \emph{Exactness of the permutation calibration.} The randomisation hypothesis:
when the joint law is invariant under permutation, a permutation test has exact level
$\alpha$. Lehmann and Romano, \emph{Testing Statistical Hypotheses}, \S 15.2; stated for
this setting in Romano and Tirlea (2022), \S\S 2--3. \PROVED
\item \emph{The max-over-lags form.} A maximum statistic must be calibrated as a maximum,
not by combining per-lag tests, or the level is not $\alpha$. This is elementary
multiplicity control and needs no special citation, but it is the reason the
implementation calibrates $\max_h$ against the permutation distribution of $\max_h$
rather than running $H$ tests.
\end{enumerate}

\paragraph{Calibration.} Within-segment permutation, upper tail, exact under
exchangeability by piece 2 above. \PROVED{} \emph{Within-segment} is load-bearing:
permuting across a read gap would manufacture inter-event gaps the carve refuses to
assert, which is the same objection that forbids concatenating segments for C3.

\paragraph{A naming caution that should be recorded.}
The word \emph{studentized} here means \emph{each lag is put on its own scale}, because
different lags have different within-segment pair counts and therefore different null
spreads. That is a lag-comparability correction and it is well motivated: at ten
segments of five the lag pair counts run from forty down to ten and the null standard
deviation ratio reaches $1.64$. \MEASURED

This is \emph{not} the studentization of Romano and Tirlea (\emph{Journal of Time Series
Analysis}, 2022; arXiv:2009.03170), which divides by a consistent estimator of the
asymptotic standard error in order to restore level $\alpha$ when the null is
\enquote{uncorrelated} rather than \enquote{independent and identically distributed}.
Their caveat is that a plain permutation test is not level $\alpha$ under an
uncorrelated-but-dependent null, with naive rejection reaching $0.233$ at $n = 50$
against a nominal $0.05$. \SIMULATED

\medskip
\noindent
The caveat does not bite here, and the reason should be stated rather than assumed: the
null this battery tests is the renewal null, which is exchangeability, not mere
uncorrelatedness. A permutation test is exact under exchangeability. But the shared word
invites a reviewer to think Romano-Tirlea has been applied when it has not, so the
documentation should say which of the two problems the code solves.

\paragraph{Censoring.} None. Complete windows only.

% ==============================================================
\section{C6 -- Permutation exchangeability, portmanteau}
% ==============================================================

\paragraph{Object.} $Q = \sum_{h=1}^{H} m_h\, r_h^2$, where $m_h$ is the number of
within-segment pairs at lag $h$.

\paragraph{Null.} Exchangeability, tested directly.

\paragraph{Directed at.} Dependence spread thinly across several lags. The $m_h$
weighting allocates the statistic's attention in proportion to how much evidence each
lag rests on, which is the opposite trade to C5's maximum. Neither substitutes for the
other.

\paragraph{Provenance.} The portmanteau form, a weighted sum of squared
autocorrelations across lags, is the rank analogue of the Box-Pierce and Ljung-Box
statistics (Box and Pierce, \emph{JASA} 65:1509--1526, 1970; Ljung and Box,
\emph{Biometrika} 65:297--303, 1978). Those use asymptotic $\chi^2_H$ critical values and
are known to be badly oversized under dependence at small $n$: Romano and Tirlea report
Ljung-Box at $0.187$ and Box-Pierce at $0.174$ against a nominal $0.05$ at $n=50$.
\SIMULATED{} The $\chi^2$ route is therefore \emph{not} taken here; only the permutation
route is.

\paragraph{Calibration.} Within-segment permutation, upper tail. Exact under
exchangeability by the randomisation hypothesis (Lehmann and Romano, \S 15.2). \PROVED

\textbf{This is the one check in the battery whose validity rests on no approximation and
no transcription.} C1 and C2 rest on an asymptotic limit and on a correctly transcribed
equation; C3 rests on an external package and a simulated null; C5 shares C6's exactness
but adds a max-statistic construction. C6 is a permutation test of exactly the null it
states. That makes it the natural reference point when checks disagree on a cell: the
burden falls on the other check.

\paragraph{Why $m_h$ weighting rather than uniform.} The number of within-segment pairs
available at lag $h$ falls as $h$ grows, and falls fastest when segments are short.
Weighting each squared autocorrelation by its pair count $m_h$ allocates the statistic's
attention in proportion to the evidence behind each term, which is the same principle
behind the $(n-h)$ weighting in Ljung-Box. Without it, a high-variance estimate from a
lag with ten pairs enters on equal terms with a stable one from forty.

\paragraph{Direction.} $Q$ is a sum of squares, therefore one-sided by construction, and
says nothing about the sign of any dependence it finds. The lag-1 rank correlation is
carried separately for that.

\paragraph{Censoring.} None. Complete windows only.

% ==============================================================
\section{C4 -- Lin-Wei-Ying, and why it is not in the battery}
\label{sec:lwy}
% ==============================================================

\paragraph{Object.} Cumulative sums of martingale-based residuals from a \emph{fitted}
Cox or Aalen regression. Lin, Wei and Ying, \emph{Biometrika} 80(3):557--572, 1993.
\PROVED

\paragraph{What it tests.} Whether the imposed regression \emph{structure} is correct:
functional form of the covariates, the link function, proportionality of hazards. Under
the null that the model is correctly specified, the cumulative residual process converges
to a zero-mean Gaussian process, and the observed path is compared against realisations
simulated with the covariates and estimated parameters held fixed.

\paragraph{What it does not test.} It does not test whether the durations are Weibull,
exponential, or anything else. It does not test independence in any sense. Fitting a
parametric family to the durations does not give it a null, because the object it checks
is a covariate structure, not a distribution family. See
Section~\ref{sec:semiparametric}.

\paragraph{Why it was in the candidate set at all.} Two circumstantial reasons, neither
of them principled. First, it handles right censoring natively, which almost nothing
else in the serial-independence set does. Second, its null is the multiplicative
intensity model, which is the same object that licenses the band and the log-rank test.
That looks like a match but is not, because the version Lin-Wei-Ying checks is the one
\emph{with covariates}, whereas the case here is the degenerate single-unit one in which
the model collapses to independent and identically distributed gaps.

\paragraph{The correct relationship.} One model, two specialisations:
%
\begin{center}
\begin{tabular}{@{}ll@{}}
\toprule
covariates plus free baseline & Cox or Aalen $\rightarrow$ Lin-Wei-Ying checks it\\
no covariates, single unit, deterministic age hazard & renewal $\rightarrow$ C1, C2, C3, C5, C6 check it\\
\bottomrule
\end{tabular}
\end{center}

\paragraph{Status.} Parked in a World C backlog. It becomes live when a regression with
covariates is fitted, for example night, threshold, qubit, or gate-voltage stability. It
is not carried as a permanently inapplicable row in the ledger, because a row that can
never be computed invites the same question at every reading.

% ==============================================================
\section{Shape statistics: does the margin predict remaining life?}
% ==============================================================

These are not licensing checks. They answer the separate question of whether the raw
signal, rather than the carved window, suffices to predict how long a window will last.
They are grouped here because the choice among them has been questioned and the reasons
should be written down.

\subsection{The four candidates and what separates them}

\begin{center}
\small
\begin{tabular}{@{}p{0.16\textwidth}p{0.08\textwidth}p{0.30\textwidth}p{0.36\textwidth}@{}}
\toprule
\textbf{Measure} & \textbf{Sym.} & \textbf{Detects} & \textbf{Null distribution}\\
\midrule
Pearson $r$ & yes & linear only; assumes normality & exact under bivariate normal\\
Spearman $\rho$ & yes & monotone, linear or not & exact permutation; asymptotic normal\\
Distance corr.\ $\mathcal{R}$ & yes & arbitrary dependence; $=0$ iff independent &
non-pivotal; permutation or bootstrap\\
Chatterjee $\xi$ & \textbf{no} & complete (functional) dependence; $=0$ iff independent,
$=1$ iff $Y = g(X)$ &
$\sqrt{n}\,\xi \to N(0, 2/5)$ under independence, \emph{tie-free case}\\
Mutual information & yes & arbitrary statistical dependence &
no closed form; estimator-dependent\\
\bottomrule
\end{tabular}
\end{center}

\paragraph{The asymmetry is the argument.} The engineering question is directional:
\emph{can remaining life be read off the margin?} Distance correlation and mutual
information are symmetric. They answer \enquote{is there dependence}, which is a weaker
and less useful statement than \enquote{is forward time a function of margin}. Chatterjee
(2021, p.~2010) makes the asymmetry deliberate:
\enquote{we may want to understand if $Y$ is a function of $X$, and not just if one of the
variables is a function of the other.} \ASSERTED

\paragraph{The symmetry problem that motivated $\xi$ in the first place.} A symmetric
excursion drives Spearman $\rho$ between a hump-shaped margin and a monotone forward time
to zero \emph{by construction}, regardless of how predictable the excursion is. Matched
simulation returned falling-limb $\rho$ of $0.74$ to $0.79$ on data whose pooled $\rho^2$
was $0.001$. \MEASURED{} So $\rho^2 \approx 0$ alone cannot support the claim that the
margin carries no information; the supportable claim is narrower, that there is no
monotone rule from margin alone.

\subsection{Reference and definition for \texorpdfstring{$\xi$}{xi}}

Chatterjee, S. (2021), \enquote{A New Coefficient of Correlation}, \emph{Journal of the
American Statistical Association} 116(536):2009--2022; preprint arXiv:1909.10140.
R package \texttt{XICOR}; a third-party Python package exists.

Sort the pairs so that $X_{(1)} \le \dots \le X_{(n)}$, breaking ties uniformly at
random. Let $r_i$ be the number of $j$ with $Y_j \le Y_{(i)}$, and $\ell_i$ the number
with $Y_j \ge Y_{(i)}$. Then
%
\begin{equation}
\xi_n(X,Y) \;=\; 1 - \frac{n \sum_{i=1}^{n-1} |r_{i+1} - r_i|}
{2 \sum_{i=1}^{n} \ell_i (n - \ell_i)} ,
\label{eq:xi}
\end{equation}
%
which in the tie-free case reduces to
$1 - 3\sum_{i=1}^{n-1}|r_{i+1}-r_i| / (n^2-1)$.
Properties: $\xi \to 0$ if and only if $X$ and $Y$ are independent, and $\xi \to 1$ if
and only if $Y$ is almost surely a measurable function of $X$. \PROVED

\subsection{Three caveats on \texorpdfstring{$\xi$}{xi} that apply directly here}

\begin{enumerate}[leftmargin=1.6em,itemsep=3pt]
\item \textbf{Ties break the $2/5$ null.} The variance $2/5$ is derived for continuous,
tie-free data. Excursion durations on a read grid are heavily tied, and at loose
thresholds most windows are a single read. Where ties dominate, the closed-form $p$-value
must not be used and a permutation calibration should replace it. This is the same
quantisation factor that threatens C3, C5 and C6, and it belongs in the same bench
column.

\item \textbf{Small-sample convergence is contested.} Chatterjee states the normal
approximation is \enquote{roughly valid even for $n$ as small as 20}. \ASSERTED{}
Independent practitioner commentary suggests $n$ of several hundred before the
approximation is comfortable. With window counts from about 35 to 355 this project sits
squarely in the disputed range, so $\xi$'s $p$-value should be permutation-calibrated in
the bench regardless of the closed form.

\item \textbf{Lower power than Spearman against smooth monotone alternatives.} $\xi$ buys
consistency against non-monotone alternatives and pays for it in power where the
relationship is in fact monotone. This is why $\xi$ and $\rho$ ship together rather than
one replacing the other. The diagnostic signature worth naming: $\rho^2$ near zero with
$\xi$ well above its null means structure exists and is not monotone.
\end{enumerate}

\subsection{Distance correlation: keep}

Székely, Rizzo and Bakirov (2007), \emph{Annals of Statistics} 35(6):2769--2794.
$\mathcal{R} = 0$ if and only if the variables are independent, for arbitrary dependence.
Symmetric, so it does not answer the directional question, but it is a legitimate
omnibus and it is what the Qubit Health Analytics work uses, which makes it the
comparable statistic when placing this framework against that one. It is already computed
in the shape statistics and should stay.

\subsection{Mutual information: declined, with the reason}

Mutual information detects arbitrary statistical dependence and is the natural
information-theoretic answer. It is declined here on three grounds, and the grounds are
practical rather than principled:

\begin{enumerate}[leftmargin=1.6em,itemsep=2pt]
\item \textbf{No closed-form null.} Every $p$-value must be permutation-calibrated, which
is affordable but removes the one advantage a closed form would confer.
\item \textbf{Estimator bias at these sample sizes.} Binning and $k$-nearest-neighbour
estimators are both biased at $n$ in the tens to low hundreds, and the bias direction
depends on the estimator and its tuning.
\item \textbf{A tuning parameter.} Bin width or $k$ must be chosen, and the estimate moves
with it. This is the same objection already used to retire the block bootstrap from
calibration routing, and applying it inconsistently would be worse than applying it
strictly.
\end{enumerate}

\noindent
If a reviewer asks why the Qubit Health Analytics quartet was reduced to a trio, that is
the answer, and it is a defensible one: the two symmetric omnibus measures are
represented by distance correlation, the monotone case by Spearman, and the directional
functional question by $\xi$, which their set does not cover at all.

% ==============================================================
\section{Validation protocol for transcribed statistical code}
\label{sec:validation}
% ==============================================================

Transcribing a published equation is not a neutral act. The protocol below is the
standing requirement for \emph{every} statistical routine in this project that is not a
direct library call, and it is stated as a general rule rather than a one-off exercise so
that a routine added later cannot escape it.

\subsection{Four tiers, all required}

\begin{description}[leftmargin=1.8em,itemsep=4pt]
\item[Tier 1, internal identity.] The transcription reproduces an independently derived
algebraic form of the same object. Example already in place: eq.~(7) at $\gh=1$ is
checked against the textbook Anderson-Darling form
$-N - N^{-1}\sum(2i-1)[\ln u_i + \ln(1-u_{N+1-i})]$ to float precision, and eq.~(16) at
$m=1$ is checked against eq.~(4).

\item[Tier 2, published reference value.] The routine reproduces a number printed in the
literature, from data also printed in the literature. This is the tier the battery was
missing and it is the one that catches transcription errors nothing else catches. Any
discrepancy is documented rather than tuned away.

\item[Tier 3, null calibration.] Under a simulated null the $p$-values are uniform, or
the empirical rejection rate matches nominal within Monte Carlo error, at the sample
sizes actually used. For exact tests this is a property that must hold and its failure is
a bug; for asymptotic tests it is a measurement of asymptotic error, not a pass or fail.

\item[Tier 4, cross-implementation.] Where an independent implementation exists, the
statistic agrees with it on a fixed input to a pinned tolerance, with the reference values
committed as fixtures rather than recomputed at test time.
\end{description}

\subsection{Tier 2 is available for C1 and C2 and should be taken}

The paper carries two worked examples with numbers, and the first carries its data.

\paragraph{Load-haul-dump machine data.} Kumar, Klefsjö and Granholm (1989), reproduced
as the paper's Table 1: 36 failure times in hours, time censored at $\tau=2000$. The
published targets are

\begin{center}
\small
\begin{tabular}{@{}llll@{}}
\toprule
\textbf{Quantity} & \textbf{Paper} & \textbf{This code} & \textbf{Status}\\
\midrule
$\hat\mu$ (complete gaps) & 54.72 & 54.72 & exact\\
$\hat\sigma$ (complete gaps) & 48.61 & 47.93 & \textbf{differs, see below}\\
$\gh$ (complete gaps) & 0.888 & 0.876 & \textbf{differs, see below}\\
$\tilde\mu$ (eq.~10) & 55.56 & 55.56 & exact\\
$\tilde\sigma$ (eq.~10) & 47.23 & 47.23 & exact\\
$\tilde\gamma$ (eq.~10) & 0.850 & 0.850 & exact\\
$\sigma^{*}$ (eq.~11) & 42.77 & 42.77 & exact\\
Laplace statistic & 0.605 & 0.605 & exact\\
LR with $\sigma^{*}$ & 0.774, $p=0.44$ & 0.774, $p=0.44$ & exact\\
LR with $\gh$ & 0.681, $p=0.50$ & 0.691, $p=0.49$ & follows from the above\\
\bottomrule
\end{tabular}
\end{center}

\paragraph{The one discrepancy, and it is the paper's, not the code's.}
The paper's Appendix 1 defines the estimator with the population divisor,
$\hat\sigma^2 = N^{-1}\sum(X_i - \hat\mu)^2$, and proves consistency in that form. Its
Table 2 reports $48.61$, which is the \emph{unbiased} $1/(N-1)$ value:
$47.93 \times \sqrt{36/35} = 48.61$ to two decimals, and
$48.61/54.72 = 0.888$, giving $0.605/0.888 = 0.681$ as printed. So the paper's text and
its table use different divisors. \MEASURED

The implementation follows the \emph{text}, with the stated reason that matching
eq.~(10)'s divisor makes the two $\gamma$ estimators differ only by the residual term.
That is a defensible choice and the resulting $1.4\%$ shift in $\gh$ moves no verdict.
But it must be recorded, because a golden test that pins $0.888$ would fail against
correct code, and a golden test that pins $0.876$ without explanation looks like a
transcription error to a reviewer holding the paper.

\paragraph{Small bowel motility data.} Aalen and Husebye (1991), 19 persons, 80 complete
periods, one censored period each. This is an $m=19$ example and therefore the only
published exercise of eq.~(16) in its multi-process form. Published targets:
$\hat\mu = 98.76$, $\hat\sigma = 52.62$, $\gh = 0.533$; Weibull fit
$\tilde\mu = 104.49$, $\tilde\sigma = 52.45$, $\tilde\gamma = 0.502$; Laplace $=1.95$;
$\mathrm{LR}^m = 1.95/0.533 = 3.67$ with $p = 0.00024$, against $p = 0.051$ for the
corresponding HPP null. The data are in Aalen and Husebye (1991), not in the paper, so
this fixture requires obtaining that dataset. It is worth the effort: nothing else
validates the $m>1$ path against a published number.

\subsection{Tier 2 for the checks that have no published example}

C5 and C6 are assembled, so no worked example exists. Their Tier 2 substitute is
threefold and each part is checkable:

\begin{enumerate}[leftmargin=1.6em,itemsep=2pt]
\item \emph{Exactness is a testable property, not an aspiration.} Under an exchangeable
null the permutation $p$-value must be uniform on the attainable grid. A
Kolmogorov-Smirnov test of the $p$-value distribution against uniform, over a few
thousand replicates, is the strongest single statement available about these two checks
and it does not depend on any citation.
\item \emph{The lag-1 statistic has an external counterpart.} \texttt{randtests::bartels.rank.test}
computes the rank von Neumann ratio, and base R's
\texttt{cor(rank(x[-n]), rank(x[-1]), method="spearman")} computes the lag-1 rank
autocorrelation directly. Both are Tier 4 cross-checks on the statistic.
\item \emph{Romano and Tirlea's Table 5.1 is a reproducible target for the bench.} At
$n=50$ under a true independent and identically distributed null they report studentized
$0.0465$ and naive $0.0493$; under $m$-dependence with $m=1$, studentized $0.0611$ and
naive $0.1388$. Their data-generating process is raw Gaussian sequences rather than
carved durations, so this is an analogue and not a fixture, but a bench that cannot
reproduce the qualitative separation is a broken bench. \SIMULATED
\end{enumerate}

\subsection{Tier 2 and 4 for \texorpdfstring{$\xi$}{xi}}

$\xi$ is currently computed in the shape statistics without any of the four tiers, which
is an inconsistency: it is transcribed code held to a lower standard than the checks. It
should be brought up to the same bar.

\begin{description}[leftmargin=1.8em,itemsep=3pt]
\item[Tier 1.] The tie-free simplification $1 - 3\sum|r_{i+1}-r_i|/(n^2-1)$ must agree
with the general form \eqref{eq:xi} on tie-free input.
\item[Tier 4.] \texttt{XICOR::xicor} is on CRAN and returns both the coefficient and the
asymptotic $p$-value. With R now available this is a direct fixture: commit a handful of
inputs and the R answers, and pin them.
\item[Tier 3, and this is the one that matters here.] The $2/5$ asymptotic variance is
derived for continuous tie-free data. Our durations are lattice-valued and heavily tied.
The required experiment is a three-way comparison on the same data: the published
asymptotic $p$-value, our closed-form $p$-value, and our permutation $p$-value, across
increasing tie fraction. Where the closed form and the permutation diverge, the closed
form is not usable and the ledger must say so.
\end{description}

\noindent
The same three-way figure is the natural validation artifact for the whole document:
one panel per instrument, showing published value, our asymptotic value, and our
permutation value, with the sample sizes and tie fractions the project actually operates
at. That figure is what licenses an instrument for the pipeline.

\subsection{Distance correlation}

Distance correlation is symmetric and consistent against arbitrary dependence, and is
what the comparable qubit-calibration literature uses, so it belongs in the shape
statistics on equal footing with $\xi$ and $\rho$. Its validation follows the same tiers:
Tier 1 against the double-centred distance matrix definition, Tier 4 against
\texttt{energy::dcor} in R or the \texttt{dcor} package in Python, Tier 3 by permutation
since the null is not pivotal. It should be rendered as a heatmap alongside $\xi$ and
$\rho$ so the three are read together rather than one standing in for the others.

% ==============================================================
\section{Code availability of the cited methods}
\label{sec:codestatus}
% ==============================================================

Where an implementation exists, use it. Where it does not, this table records what was
looked for and what the authors say, so that transcription is a documented last resort
rather than a default.

\begin{center}
\small
\setlength{\tabcolsep}{4pt}
\begin{tabular}{@{}>{\raggedright\arraybackslash}p{0.20\textwidth}
                 >{\raggedright\arraybackslash}p{0.17\textwidth}
                 >{\raggedright\arraybackslash}p{0.24\textwidth}
                 >{\raggedright\arraybackslash}p{0.30\textwidth}@{}}
\toprule
\textbf{Source} & \textbf{Needed for} & \textbf{Code status} & \textbf{Action}\\
\midrule
Kvaløy and Lindqvist (2020), \emph{Technometrics} 62(1) &
C1, C2, and CvM, ELR, KS if added &
\textbf{On request.} Their \S 7 closing line: \enquote{R-code for the tests can be
obtained from the authors.} No CRAN, PyPI or public repository located. &
\textbf{Write to the authors.} Correspondence address on the Technometrics record is
\texttt{jan.t.kvaloy@uis.no} (University of Stavanger); Bo Henry Lindqvist is at NTNU
Trondheim. Ask for the code behind Table 3, which would also supply reference values for
AD, CvM, KS and ELR on the load-haul-dump data.\\[3pt]

Lindqvist and Kvaløy (2024), \emph{ASMB}, doi:10.1002/asmb.2848 &
integrated and supremum variants, a possible future extension &
Not located. Same group. &
Ask in the same message. They report simulated critical values for ICvM, IKS and SELR1
which would be reference values if those tests are ever added.\\[3pt]

Marsaglia and Marsaglia (2004), \emph{JSS} 9(2) &
C2 limiting CDF &
\textbf{Public with the article.} \emph{Journal of Statistical Software} publishes source
alongside every paper. &
Download and use as a Tier 4 fixture rather than contacting anyone.\\[3pt]

Lawless, Çiğşar and Cook (2012), \emph{Technometrics} 54 &
validity of permutation under time censoring &
Not located. &
Low priority: the result is cited, not computed, so no code is needed to use it.\\[3pt]

Romano and Tirlea (2022), \emph{JTSA} &
C5 studentization, size targets &
No package located; simulation tables are in the paper. &
Use the printed Table 5.1 numbers as bench targets. Contact only if the exact generator
is needed.\\[3pt]

Chatterjee (2021), \emph{JASA} 116(536) &
$\xi$ &
\textbf{Public tool.} \texttt{XICOR} on CRAN, with an independence test. A third-party
Python package exists. &
Nothing to request. Use as a Tier 4 fixture.\\[3pt]

Genest-Rémillard, Ghoudi-Kulperger-Rémillard, Kojadinovic-Yan &
C3 &
\textbf{Public tool.} \texttt{copula} on CRAN. &
Nothing to request. Already the implementation.\\[3pt]

Székely, Rizzo and Bakirov (2007) &
distance correlation &
\textbf{Public tool.} \texttt{energy} on CRAN; \texttt{dcor} on PyPI. &
Nothing to request.\\[3pt]

Lin, Wei and Ying (1993) &
C4, parked &
\textbf{Public tool.} \texttt{timereg::cum.residuals}, \texttt{gof::cumres.coxph}. &
Nothing to request until a regression is fitted.\\
\bottomrule
\end{tabular}
\end{center}

\medskip
\noindent
Only one row requires a message, and it is the row that matters most: the two tests the
battery leans on hardest are the two with no public implementation in any language. A
short, specific request naming the paper, the table, and the intended use is more likely
to succeed than a general one, and reference values for Table 3 would upgrade C2 from
Tier 1 to Tier 2 in a single reply.

% ==============================================================
\section{Implementation provenance}
% ==============================================================

Every statistical routine in the battery, with what it calls and where the mathematics
comes from.

\begin{center}
\small
\begin{tabular}{@{}p{0.10\textwidth}p{0.22\textwidth}p{0.28\textwidth}p{0.32\textwidth}@{}}
\toprule
\textbf{Check} & \textbf{Implementation} & \textbf{Source of the mathematics} &
\textbf{Why not a library}\\
\midrule
C1 & transcribed, pure Python; \texttt{scipy.stats.norm.sf} for the normal tail &
Kvaløy-Lindqvist eq.~(16), collapsing to eq.~(4) at $m=1$ &
No CRAN or PyPI package implements the time-censored form. The paper's own closing line:
\enquote{R-code for the tests can be obtained from the authors.}\\[3pt]
C2 & transcribed, pure Python; Marsaglia-Marsaglia \texttt{adinf} for the limiting CDF &
Kvaløy-Lindqvist eq.~(7); Marsaglia and Marsaglia, \emph{JSS} 9(2), 2004 &
Same. \texttt{scipy.stats.anderson} is a goodness-of-fit test for named distributions and
is a different statistic.\\[3pt]
C3 & \textbf{imported} over a file bridge to \texttt{Rscript} &
\texttt{copula::serialIndepTest}, \texttt{copula::serialIndepTestSim} &
Library exists and is used. No maintained Python equivalent.\\[3pt]
C5 & assembled from primitives; \texttt{scipy.stats.rankdata} &
Rank autocorrelation with permutation calibration; Romano-Tirlea for the level caveat &
No canonical package for permutation rank autocorrelation on segmented series.
\texttt{pastecs::trend.test} is a trend test, not a serial test, and is not a
substitute.\\[3pt]
C6 & assembled from primitives; shares C5's harness &
Portmanteau rank statistic, exact permutation &
Same.\\[3pt]
C4 & \textbf{not implemented} &
Lin, Wei and Ying, \emph{Biometrika} 80, 1993 &
\texttt{timereg::cum.residuals} and \texttt{gof::cumres.coxph} both exist, but there is
no fitted regression for them to check.\\
\bottomrule
\end{tabular}
\end{center}

\paragraph{Summary of the audit.} Of five implemented checks, one imports a library and
four transcribe published equations. The four are not reinvention: no package implements
the time-censored Kvaløy-Lindqvist class in any language, and the authors distribute code
on request rather than through a repository. For C5 and C6 the statistic is elementary and
the difficulty is entirely in the segmented permutation harness, which no package
provides because no package knows about read gaps.

\paragraph{Where a library should be preferred and currently is not.} If the
Cramér-von Mises test of Section~\ref{sec:missing} is added, it should reuse the same
transcription path, since it comes from the same paper and no package implements it
either. For $\xi$, the R package \texttt{XICOR} exists and includes the independence test;
if a shape statistic ever needs an authoritative reference value, that package is the
cross-check, in the same way R-generated golden fixtures should pin C3.

% ==============================================================
\section{Consolidated table}
% ==============================================================

\begin{center}
\footnotesize
\setlength{\tabcolsep}{4pt}
\begin{tabular}{@{}>{\raggedright\arraybackslash}p{0.05\textwidth}
                 >{\raggedright\arraybackslash}p{0.13\textwidth}
                 >{\raggedright\arraybackslash}p{0.20\textwidth}
                 >{\raggedright\arraybackslash}p{0.20\textwidth}
                 >{\raggedright\arraybackslash}p{0.08\textwidth}
                 >{\raggedright\arraybackslash}p{0.21\textwidth}@{}}
\toprule
\textbf{Id} & \textbf{Null} & \textbf{Directed at} & \textbf{Calibration} &
\textbf{Cens.} & \textbf{Status}\\
\midrule
C1 & renewal & monotone trend & asymptotic $N(0,1)$; permutation & native & implemented\\
C2 & renewal & monotone and non-monotone trend & limiting AD; permutation & native & implemented\\
C3 & serial independence & any serial dependence & simulated null (R) & none & unassessed\\
C5 & exchangeability & localised neighbour dep. & permutation & none & implemented\\
C6 & exchangeability & spread neighbour dep. & permutation & none & implemented\\
\midrule
CvM & renewal & monotone and non-monotone trend & limiting CvM; permutation & native & \textbf{recommended}\\
ELR & renewal & symmetric non-monotone & asymptotic $N(0,1)$ & native & declined, tuning parameter\\
KS & renewal & monotone trend & Kolmogorov & native & declined, dominated\\
Mann & renewal & monotone trend & rank & limited & declined, redundant with C1\\
C4 & fitted regression & misspecification & simulated Gaussian null & native & parked, World C\\
\bottomrule
\end{tabular}
\end{center}

% ==============================================================
\section*{References}
% ==============================================================
\addcontentsline{toc}{section}{References}

\begin{description}[leftmargin=1.2em,labelindent=0em,itemsep=1pt,font=\normalfont]
\item Andersen, P.~K., Borgan, Ø., Gill, R.~D. and Keiding, N. (1993).
\emph{Statistical Models Based on Counting Processes}. Springer.
\item Anderson, T.~W. and Darling, D.~A. (1952). Asymptotic theory of certain goodness of
fit criteria based on stochastic processes. \emph{Annals of Mathematical Statistics}
23:193--212.
\item Aalen, O.~O. and Husebye, E. (1991). Statistical analysis of repeated events
forming renewal processes. \emph{Statistics in Medicine} 10:1227--1240.
\item Bartels, R. (1982). The rank version of von Neumann's ratio test for randomness.
\emph{Journal of the American Statistical Association} 77(377):40--46.
\item Billingsley, P. (1999). \emph{Convergence of Probability Measures}, 2nd ed. Wiley.
\item Box, G.~E.~P. and Pierce, D.~A. (1970). Distribution of residual autocorrelations
in autoregressive-integrated moving average time series models. \emph{Journal of the
American Statistical Association} 65:1509--1526.
\item Chatterjee, S. (2021). A new coefficient of correlation. \emph{Journal of the
American Statistical Association} 116(536):2009--2022.
\item Genest, C. and Rémillard, B. (2004). Tests of independence and randomness based on
the empirical copula process. \emph{TEST} 13(2):335--369.
\item Ghoudi, K., Kulperger, R. and Rémillard, B. (2001). A nonparametric test of serial
independence for time series and residuals. \emph{Journal of Multivariate Analysis}
79(2):191--218.
\item Kojadinovic, I. and Yan, J. (2011). Tests of serial independence for continuous
multivariate time series based on a Möbius decomposition of the independence empirical
copula process. \emph{Annals of the Institute of Statistical Mathematics}
63(2):347--373.
\item Kumar, U., Klefsjö, B. and Granholm, S. (1989). Reliability investigation for a
fleet of load haul dump machines in a Swedish mine. \emph{Reliability Engineering and
System Safety} 24:341--361.
\item Kvaløy, J.~T. and Lindqvist, B.~H. (1998). TTT-based tests for trend in repairable
systems data. \emph{Reliability Engineering and System Safety} 60:13--28.
\item Kvaløy, J.~T. and Lindqvist, B.~H. (2020). A class of tests for trend in time
censored recurrent event data. \emph{Technometrics} 62(1):101--115.
Preprint arXiv:1802.08339.
\item Lawless, J.~F., Çiğşar, C. and Cook, R.~J. (2012). Testing for monotone trend in
recurrent event processes. \emph{Technometrics} 54:147--158.
\item Lewis, P.~A.~W. and Robinson, D.~W. (1974). Testing for a monotone trend in a
modulated renewal process. In \emph{Reliability and Biometry}, SIAM, pp.~163--182.
\item Lehmann, E.~L. and Romano, J.~P. (2005). \emph{Testing Statistical Hypotheses},
3rd ed. Springer. (\S 15.2, randomisation tests.)
\item Lin, D.~Y., Wei, L.~J. and Ying, Z. (1993). Checking the Cox model with cumulative
sums of martingale-based residuals. \emph{Biometrika} 80(3):557--572.
\item Lindqvist, B.~H., Elvebakk, G. and Heggland, K. (2003). The trend-renewal process
for statistical analysis of repairable systems. \emph{Technometrics} 45:31--44.
\item Lindqvist, B.~H. and Kvaløy, J.~T. (2024). New tests for trend in time censored
recurrent event data. \emph{Applied Stochastic Models in Business and Industry}.
doi:10.1002/asmb.2848.
\item Mann, H.~B. (1945). Nonparametric tests against trend. \emph{Econometrica}
13:245--259.
\item Ljung, G.~M. and Box, G.~E.~P. (1978). On a measure of lack of fit in time series
models. \emph{Biometrika} 65:297--303.
\item Marsaglia, G. and Marsaglia, J. (2004). Evaluating the Anderson-Darling
distribution. \emph{Journal of Statistical Software} 9(2).
\item Romano, J.~P. and Tirlea, M.~A. (2022). Permutation testing for dependence in time
series. \emph{Journal of Time Series Analysis}. Preprint arXiv:2009.03170.
\item Székely, G.~J., Rizzo, M.~L. and Bakirov, N.~K. (2007). Measuring and testing
dependence by correlation of distances. \emph{Annals of Statistics} 35(6):2769--2794.
\end{description}

\end{document}